% Method section for an ISCA-style two-column paper.
% Expected packages: amsmath, amssymb, graphicx.
% The section is self-contained and uses no algorithm/TikZ package.
% Suggested bibliography keys to provide in the main paper:
%   riscv-rvwmo, boom, assume-guarantee, cegar, ic3

\providecommand{\sys}{\ensuremath{\mu}\textsc{MCM}}
\providecommand{\traces}{\mathcal{T}}
\providecommand{\proj}{\mathop{\downarrow}}

\section{Methodology}
\label{sec:method}

\sys{} searches a detailed model of an out-of-order RISC-V memory hierarchy
for violations of its architectural memory model; once model fidelity is
established, the result lifts to the RTL implementation.  The central obstacle is scale.
An LSU, private cache, and shared cache each contain queues, arbiters, and
protocol state whose synchronous product is prohibitively large.  Merely
placing an interface between these RTL blocks does not solve the problem:
memory-ordering bugs often arise from interactions among structures within a
block, while an unconstrained interface can hide precisely those interactions.

Our key idea is to replace monolithic state exploration with a hierarchy of
\emph{event contracts}.  We partition the memory hierarchy into fine-grained
semantic components, summarize each component by the events visible at its
ports, and prove that each summary conservatively covers its implementation.
The verifier composes summaries, rather than private state machines, and
refines them only when a candidate violation requires additional internal
detail.  Architectural relations such as reads-from and coherence order are
encoded symbolically; the verifier never enumerates schedules or candidate
orders.  Under explicit contract-compatibility conditions, the dominant cost
is the sum of local analyses plus a substantially smaller interface problem,
rather than the product of all component state spaces.

Figure~\ref{fig:overview} gives an overview.  A verification task begins with
an architectural program skeleton and a target property.  Local module
specifications complete the relevant microarchitectural events.  Contract
composition connects these events across instance-scoped ports and produces a
global interface execution.  A refinement map then constructs an RVWMO
execution graph.  A satisfying assignment to the negated RVWMO property is a
candidate bug; an unsatisfiable query proves the absence of that violation
within the stated resource and time bounds.

\begin{figure*}[t]
  \centering
  \small
  \setlength{\tabcolsep}{3pt}
  \begin{tabular}{c@{\quad$\Longrightarrow$\quad}c@{\quad$\Longrightarrow$\quad}c@{\quad$\Longrightarrow$\quad}c}
    \fbox{\parbox[c][0.58in][c]{0.18\textwidth}{\centering
      Architectural seed\\and target property}}
    &
    \fbox{\parbox[c][0.58in][c]{0.18\textwidth}{\centering
      Fine-grained\\event contracts}}
    &
    \fbox{\parbox[c][0.58in][c]{0.20\textwidth}{\centering
      Hierarchical symbolic\\composition}}
    &
    \fbox{\parbox[c][0.58in][c]{0.18\textwidth}{\centering
      RVWMO graph\\and violation query}}
  \end{tabular}
  \vspace{3pt}

  \begin{tabular}{c@{\hspace{0.6in}}c@{\hspace{0.6in}}c}
    source-anchored specification
    & counterexample-guided refinement
    & RTL witness concretization
  \end{tabular}
  \caption{\sys{} verification flow.  Solid arrows transform symbolic
  executions; feedback refines only the components involved in a spurious
  counterexample.  RTL is not flattened into the global query.}
  \label{fig:overview}
\end{figure*}

\subsection{Verification Objective}
\label{sec:objective}

\paragraph{Executions and observations.}
We model a hardware execution as a finite, partially ordered set of typed
events.  An event $e=\langle\tau,\iota,\nu\rangle$ has an event type $\tau$, a
stable identity $\iota$, and a valuation $\nu$ of typed fields.  For example,
a load-request event may carry a hart, instruction identifier, queue entry,
address, byte mask, and transaction identifier.  Event identities are
instance scoped: two identically typed events at different ports are distinct
unless a connection explicitly equates them.  The partial order contains only
semantically required edges, such as pipeline succession, queue age, request
to response, and commit order.  Independent events remain unordered, avoiding
an enumeration of cycle-by-cycle interleavings.

Let $D$ be a detailed, manually constructed microarchitectural specification
and let $\traces(D)$ denote its event traces.  A projection $\pi_I$ retains the
fields and events in interface $I$ while existentially quantifying private
events and state.  Projection is therefore a semantic operation:
\begin{equation}
  \pi_I(\traces(M)) =
  \{t_I \mid \exists t\in\traces(M).\;\pi_I(t)=t_I\},
  \label{eq:semantic-projection}
\end{equation}
not a syntactic pass that deletes every constraint mentioning a hidden event.
This distinction is essential: a private replay or invalidation may disappear
from the interface trace, but the ordering constraint it induces on public
events must survive projection.

\paragraph{Two refinement obligations.}
Our first-stage method separates local abstraction from architectural
correctness.  For every detailed component specification $D_i$ and its
contract summary $C_i$, we require
\begin{equation}
  \pi_{I_i}(\traces(D_i)) \subseteq \traces(C_i).
  \label{eq:summary-refinement}
\end{equation}
Thus $C_i$ is a conservative over-approximation: it cannot omit an interface
behavior of the detailed model.  Over-approximation is required both for bug
completeness and for proving model-level safety, although it may introduce
spurious counterexamples.
At the architectural boundary, a composed microarchitectural execution is
mapped by $\alpha$ to an execution graph containing program order ($po$),
reads-from ($rf$), coherence order ($co$), from-read ($fr$), dependencies,
fences, and atomic pairs.  Correctness requires
\begin{equation}
  \alpha\!\left(\traces(\parallel_i D_i)\right)
  \subseteq \traces(\mathrm{RVWMO}).
  \label{eq:arch-refinement}
\end{equation}
Equation~\ref{eq:summary-refinement} ensures that hierarchy itself does not
hide a detailed-model behavior, while Equation~\ref{eq:arch-refinement} states
the memory-model verification target.

This stage deliberately treats fidelity between RTL component $R_i$ and
$D_i$ as an explicit assumption.  Establishing
$\pi(\traces(R_i))\subseteq\traces(D_i)$ in the presence of hand- or
LLM-written specifications is the separate model-validation problem addressed
by our second stage.  We therefore distinguish model-complete results from
RTL-complete claims: a counterexample is confirmed against RTL, whereas an
absence result applies to the detailed model until the fidelity assumption is
discharged.  This separation prevents source annotations or regression tests
from being mistaken for a proof of model completeness.

\paragraph{Scope of completeness.}
No finite method can establish the absence of every conceivable hardware bug.
\sys{} instead provides a precise bounded claim.  A task fixes the number of
harts, architectural memory operations, abstract addresses, and instances of
bounded resources such as LDQ, STQ, and MSHR entries.  Within these bounds,
\sys{} searches all executions admitted by the conservative specifications and
all architectural relations admitted by RVWMO.  We say that a bug class is
covered when its violation predicate is expressible over the architectural and
interface events.  Our current classes include ordering, forwarding, replay,
exception/kill, atomicity, dirty-data preservation, and coherence-permission
violations.  Liveness requires a separately stated fairness and progress
contract and is not inferred from a finite trace.

\subsection{Fine-Grained Event Specifications}
\label{sec:specification}

\paragraph{Semantic rather than structural partitioning.}
RTL module boundaries are too coarse for compositional memory verification.
For example, BOOM's LSU combines load and store queues, translation, replay,
forwarding, dependence checks, exceptions, and commit interaction.  Flattening
the LSU into one state machine recreates the state explosion that composition
is intended to avoid.  We instead partition at \emph{semantic cut points}:
each component owns one resource or one phase of a memory transaction and
communicates only through typed events.

For BOOM, the LSU is decomposed into load issue, LDQ-entry lifetime, store
address/data handling, STQ-entry lifetime, translation/retry, forwarding,
load--load ordering checks, store--load ordering checks, exception/kill, and
commit adapters.  The L1 is decomposed into request arbitration, tag/meta
lookup, hit handling, refill, writeback, probe handling, replacement, and an
MSHR-entry template.  The shared cache is decomposed into directory-entry,
Acquire, Probe, Release, Grant, and transaction-entry components.  This
decomposition is not hard-coded to BOOM: components implement protocol roles,
allowing the same contract vocabulary to be instantiated for a differently
organized design such as XiangShan.

\paragraph{Module semantics.}
A detailed specification module is a tuple
\begin{equation}
  D_i=\langle X_i,E_i^{\mathrm{in}},E_i^{\mathrm{out}},
  E_i^{\mathrm{priv}},Init_i,T_i,F_i\rangle,
\end{equation}
where $X_i$ is private symbolic state, $E_i^{\mathrm{in}}$ and
$E_i^{\mathrm{out}}$ are typed port events, $E_i^{\mathrm{priv}}$ are internal
events, $Init_i$ is the initial-state predicate, $T_i$ is the event/state
transition relation, and $F_i$ is a set of fairness clauses used only for
progress properties.  $T_i$ is relational rather than functional, permitting
the specification to conservatively represent arbitration, cache replacement,
and variable response latency.

Inputs may be constrained only by the module's assumption; outputs and private
state are owned by the module.  A top-level task can provide architectural
instructions, initial memory, page-table configuration, and external protocol
responses, but it cannot inject a child-private event such as an LSU replay or
a cache hit.  Those events must be derived by $T_i$.  This rule is important
for bug discovery: otherwise a trace merely describes the internal path of a
known bug instead of asking whether the implementation can produce it.

\paragraph{Parameterized resource templates.}
Queue entries, cache lines, and MSHRs are modeled by parameterized templates,
not by copying a monolithic specification for every physical entry.  A template
is instantiated only for an identity that lies in the cone of influence of the
target property.  Identities are symbolic and obey allocation, age, ownership,
and uniqueness constraints.  Symmetric identities are interchangeable until
an observable event distinguishes them; the solver therefore chooses a
canonical representative.  Entries outside the cone are not simply removed:
an aggregate interference summary conservatively represents their occupancy,
backpressure, eviction, and arbitration effects.  This combines symmetry
reduction with data independence without assuming that inactive entries cannot
delay or perturb an active transaction, and avoids multiplying the execution
by every physical entry in a large configuration.

\paragraph{Transaction-centric interfaces.}
Interfaces carry the minimum information required to preserve memory-model
behavior: operation and transaction identities, hart, program position,
aligned block address, byte mask, value or value version, queue age,
permission, and kill/commit status.  An interface does not export pipeline
registers merely because they exist in RTL.  Conversely, it must export any
fact that can alter architectural order or value.  For example, an L1-to-LSU
response exports the responding transaction and data version, while an
invalidation exports the block and version it supersedes.  These fields allow
the verifier to determine whether a completed load became stale without
exposing the entire cache state.

\subsection{Contracts and Hierarchical Composition}
\label{sec:contracts}

\paragraph{Event contracts.}
Each module exports a contract $C_i=\langle A_i,G_i,H_i\rangle$.  The
assumption $A_i$ constrains legal input histories, the guarantee $G_i$ relates
input and output histories, and the summary relation $H_i$ captures persistent
interface state, such as the current owner and version of a cache block.  A
module discharges the local proof obligation
\begin{equation}
  D_i \models A_i \Rightarrow (G_i \land H_i).
  \label{eq:local-contract}
\end{equation}
Unlike a collection of point-to-point event maps, $G_i$ can relate multiple
events across time.  It can state, for example, that a granted transaction was
previously acquired, that a killed request never produces an architectural
completion, or that a dirty owner supplies the value version returned to a
new reader.

Ports are connected by explicit, instance-scoped channels.  A connection
unifies the producer's output with the consumer's input and checks type,
direction, identity, and multiplicity.  Thus two MSHR instances may emit the
same event type without becoming implicitly connected.  One-to-one,
one-to-many, arbitration, and lossy observation channels have distinct
semantics; an ``exact'' connection denotes cardinality, rather than merely the
existence of some matching event.

\paragraph{Compatibility and composition.}
For a set of connected children, let $Env_i$ denote guarantees of the other
children and the legal top-level environment projected onto $i$'s inputs.
The children are compatible when
\begin{equation}
  Env_i \land \bigwedge_{j\neq i}G_j \Rightarrow A_i
  \quad\text{for every }i.
  \label{eq:compatibility}
\end{equation}
The parent summary is obtained by conjoining child guarantees, equating
connected channels, and existentially eliminating child-private interface
variables:
\begin{equation}
 H_P = \exists V_{\mathrm{hidden}}.\;
       Env_P \land \bigwedge_i(G_i\land H_i) \land Connect.
 \label{eq:summary}
\end{equation}
Equation~\ref{eq:summary} preserves constraints induced by hidden behavior; it
does not first choose a concrete child trace and then erase private events.
The resulting $H_P$ is itself an event contract and can be composed at the next
level, yielding an arbitrary-depth hierarchy.

Memory protocols contain cyclic dependencies---for example, an L1 waits for a
Grant while the L2 may wait for a ProbeAck.  We handle these cycles by placing
the communicating components in a strongly connected contract group.  The
group's mutually recursive guarantees are solved simultaneously and then
summarized at the group boundary.  A well-founded causality condition requires
every response chain to originate in an environment request or in retained
initial state; request/response identities prevent an unrelated transaction
from satisfying the cycle.  In practice, the solver initially exposes only
the boundary clauses and adds an internal protocol clause when a candidate
trace uses the corresponding behavior.

\paragraph{Why composition scales.}
Let $Q_i$ be the cost of checking component $i$ and $Q_I$ the cost of solving
the retained interface variables.  Flattening constructs a problem over the
synchronous product $\prod_i|X_i|$.  Contract composition instead costs
approximately
\begin{equation}
  Q_{\mathrm{hier}} \approx \sum_i Q_i + Q_I.
  \label{eq:additive-cost}
\end{equation}
This is not a complexity claim for arbitrary circuits.  Equation
~\ref{eq:additive-cost} holds when local proofs are reusable, interface width
is bounded by the active transactions, and the target property does not force
all private state into one contract group.  \sys{} actively enforces these
conditions through semantic partitioning, cone-of-influence slicing, and
counterexample-guided disclosure of private facts.  If an interface grows with
the full private state, the tool reports the lost abstraction boundary rather
than silently constructing the global product.

\subsection{Symbolic Architectural Refinement}
\label{sec:symbolic}

\paragraph{Architectural projection.}
Commit events define architectural memory operations.  A committed load emits
an architectural read only if it survives branch, exception, and memory-order
recovery; a committed store becomes a write independently of when it drains;
an AMO emits a paired read/write; and successful LR/SC emits an atomic pair.
The projection retains byte ranges, dependencies, fence modes, and
acquire/release annotations required by RVWMO.  Microarchitectural provenance,
such as a cache response version or forwarded STQ entry, constrains the
architectural $rf$ edge but does not replace the architectural operation.

The bug query is constructed directly, rather than completing an arbitrary
microarchitectural trace before checking it:
\begin{equation}
 \exists t_\mu,g.\;
   Compose(t_\mu) \land g=\alpha(t_\mu)
   \land \neg RVWMO(g).
 \label{eq:bug-query}
\end{equation}
Consequently, if both a legal and an illegal internal completion exist, the
solver searches for the illegal one.  Unsatisfiability of
Equation~\ref{eq:bug-query} establishes bounded safety of the detailed model
because its contract summaries conservatively include the detailed behaviors.

\paragraph{Relations without enumeration.}
\sys{} represents architectural choices by symbolic variables.  Each read
$r$ has a selector $src(r)$ ranging over overlapping writes and an initial
write; $rf(w,r)$ holds exactly when $src(r)=w$.  Each write $w$ has a symbolic
coherence rank $rank(w)$, with distinct ranks for writes in the same coherence
domain.  The relation $co(w_1,w_2)$ is encoded as
$rank(w_1)<rank(w_2)$.  Program order and dependency edges are obtained from
operation identities, and $fr=rf^{-1};co$ is constrained relationally.  The
encoding size is polynomial in the number of retained operations and avoids
enumerating the factorial set of coherence orders or the Cartesian product of
read-from candidates.

The same principle applies below the architectural boundary.  Event existence
is Boolean, latency is constrained by symbolic order or bounded timestamp
variables, queue allocation is represented by symbolic identities, and
arbitration selects one enabled request.  Independent events are represented
by one partial-order execution rather than all total schedules.  The solver
adds a total order only where an axiom, a shared resource, or an observed value
requires one.

\paragraph{Property-directed refinement.}
A coarse contract may admit an architectural violation that no implementation
trace realizes.  \sys{} handles this expected consequence of conservative
abstraction with counterexample-guided refinement.  Given an abstract witness,
the tool checks each participating component against its detailed event
specification, while fixing only the witness's interface events.  A feasible
check yields a concrete microarchitectural witness.  An infeasible check
returns an interface reason---for example, a missing queue-age condition or an
impossible permission transition---which is generalized into a blocking
clause and added to the parent summary.  Components outside the witness cone
are not reopened.

The process is summarized below.  Importantly, refinement adds semantic
constraints; it never resolves a spurious witness by manually choosing hit,
miss, replay, or protocol outcomes in the input trace.

\begin{figure}[t]
  \small
  \hrule
  \vspace{3pt}
  \textbf{Hierarchical violation search}
  \begin{enumerate}
    \setlength{\itemsep}{0pt}
    \setlength{\parsep}{0pt}
    \item Slice contracts using the target property and architectural seed.
    \item Solve Equation~\ref{eq:bug-query} over the current summaries.
    \item If UNSAT, report bounded safety and the active assumptions.
    \item Otherwise, concretize the candidate in participating children.
    \item If feasible, emit a microarchitectural and architectural witness.
    \item If infeasible, learn an interface clause and return to Step~2.
  \end{enumerate}
  \vspace{-2pt}
  \hrule
  \caption{Counterexample-guided hierarchical verification.}
  \label{fig:search}
\end{figure}

\subsection{Instantiation for BOOM}
\label{sec:boom-instantiation}

\paragraph{End-to-end event path.}
We instantiate the method across BOOM's LSU and L1 data cache and the shared
inclusive L2.  The top-level task contains decoded memory instructions,
initial memory, architectural dependencies, and bounded environment events.
It contains no cache hit/miss, queue index, MSHR identity, replay decision, or
coherence outcome.  The relevant leaf contracts derive the following path:
instruction allocation and translation; load/store issue, forwarding, and
replay; L1 lookup and miss handling; TileLink Acquire/Probe/Release/Grant;
dirty-data and permission transfer; response, exception recovery, and ROB
commit; and finally architectural load/store/AMO events.

A small set of cross-cutting identities links this path.  The operation
identifier remains stable from decode to commit.  A transaction identifier
tracks each cache request and response.  The aligned block address and byte
mask define overlap.  A monotonically constrained data version tracks which
write a response observes, avoiding an enumeration of concrete data values.
Queue age and commit status are exposed only at contracts that can change
ordering or recovery.  These identities are sufficient to connect LSU, L1,
and L2 behavior without exposing their complete private state.

\paragraph{Running example.}
Consider a known class of load--load violation.  An older load is delayed while
a younger same-address load completes.  A coherence action replaces the line,
after which the older load observes a newer value.  Correct recovery must
invalidate or replay the younger load before it commits.  In \sys{}, separate
contracts describe LDQ age, successful completion, coherence-version change,
ordering-failure detection, replay, and commit.  The global violation asks
whether both loads can commit with values that induce an $fr\rightarrow
rfe\rightarrow ppo$ cycle.  The query does not prescribe a TLB miss, retry, or
conflict.  If such internal events are necessary, the LSU specification must
derive them; if the coarse LSU summary admits an impossible sequence, only the
LSU contracts are refined.  This example exercises a cross-cutting property
without flattening the complete LSU and cache state.

\paragraph{Bug classes and output.}
The same workflow targets store-to-load forwarding errors, lost ordering
replays, committed stores removed by exception recovery, fence/AMO ordering,
LR/SC atomicity, stale refill data, lost dirty writeback, and incorrect
coherence ownership.  Every reported violation contains two linked
certificates: a microarchitectural event witness explaining the internal path,
and an architectural graph explaining the violated RVWMO axiom.  Source
anchors associate each leaf rule with the RTL expressions and revision from
which it was constructed.  These artifacts make a result auditable and form
the interface to later RTL-conformance validation and automated specification
generation; neither is placed in the trusted core of the first-stage method.

\paragraph{Soundness argument.}
For completeness, suppose every leaf satisfies
Equation~\ref{eq:summary-refinement}, every composition discharges
Equation~\ref{eq:compatibility}, and summaries are constructed by
Equation~\ref{eq:summary}.  Contract composition preserves trace inclusion by
existential elimination, so the composition of detailed models is included in
the root summary.  If Equation~\ref{eq:bug-query} is unsatisfiable at the root,
no bounded trace of the detailed model can project to an RVWMO violation.
Conversely, a feasible candidate that passes leaf concretization denotes an
execution of the detailed microarchitectural specifications and an explicit
architectural violation.  Extending the safety result to RTL additionally
requires $\pi(\traces(R_i))\subseteq\traces(D_i)$ for every leaf.  This
argument separates four outcomes that are often conflated: model-level bounded
safety, conditional RTL safety, a concretized bug, and a spurious abstract
witness that requires refinement.
